% Môn: Toán
% Lớp: 10
% Chương: Chương I. Mệnh đề và tập hợp
% Bài: Bài 2. Tập hợp và các phép toán trên tập hợp · Các tập số, khoảng, đoạn và nửa khoảng
% Số câu: 56
% App đọc file này trực tiếp — sửa rồi Commit trên GitHub, bấm Đồng bộ GitHub .tex

% Bài 2. Tập hợp và các phép toán trên tập hợp



% ===== Câu 8 =====
% ID: T10C1B2-02-TN
% Mức: NB
\dangbt{Các tập số, khoảng, đoạn và nửa khoảng}
\begin{ex}
% ID: T10C1B2-02-TN
% Mức: NB
Cách viết nào sau đây là $\textbf{sai}$?
\choice
{$(0 ;+\infty)=\{x \in \mathbb{R}| x>0\}$}
{$[0 ; 5]=\{x \in \mathbb{R}| 0 \leqslant x \leqslant 5\}$}
{$(0 ; 5)=\{x \in \mathbb{R}| 0\lt x<5\}$}
{\True $[-1 ; 5)=\{-1 ; 0 ; 1 ; 2 ; 3 ; 4\}$}
\loigiai{
Vì $[-1 ; 5)=\{x \in \mathbb{R} \mid -1 \leqslant x<5\}$ nên cách viết sau là sai: «$[-1 ; 5)=\{-1 ; 0 ; 1 ; 2 ; 3 ; 4\}$».
}
\end{ex}

% ===== Câu 38 =====
% ID: T10C1B2-20-TN
% Mức: NB
\begin{ex}
% ID: T10C1B2-20-TN
% Mức: NB
Cho tập hợp $A=[1;5]$, phát biểu nào sau đây đúng?
\choice
{$A=\{x \in \mathbb{R} \mid 1 \lt  x \lt  5\}$}
{\True $A=\{x \in \mathbb{R} \mid 1 \leq x \leq 5\}$}
{$A=\{x \in \mathbb{R} \mid 1 \leq x \lt  5\}$}
{$A=\{x \in \mathbb{R} \mid 1 \lt  x \leq 5\}$}
\loigiai{
Ta có $A=[1;5]=\{x \in \mathbb{R} \mid 1 \leq x \leq 5\}$. Đoạn $[1;5]$ bao gồm cả hai đầu mút $1$ và $5$, tức là $x$ thỏa mãn $1 \leq x \leq 5$. Vậy phát biểu đúng là phương án B.
}
\end{ex}

% ===== Câu 42 =====
% ID: T10C1B2-24-TN
% Mức: TH
\begin{ex}
% ID: T10C1B2-24-TN
% Mức: TH
Cho $A=\{1;5\};B=\{1;3;5\}$. Chọn kết quả đúng trong các kết quả sau
\choice
{$A \cup B = \{1;3;5\}$}
{$A \cap B = \{1\}$}
{\True $A \subset B$}
{$B \subset A$}
\loigiai{
Ta có $A=\{1;5\}$ và $B=\{1;3;5\}$.
Kiểm tra các phương án:
$A \cup B = \{1;3;5\}$ (Sai)
$A \cap B = \{1;5\}$ (Sai)
$A \subset B$ vì mọi phần tử của $A$ đều thuộc $B$ (Đúng)
$B \subset A$ (Sai)
Vậy phương án $C$ là đúng.
}
\end{ex}
